% This is samplepaper.tex, a sample chapter demonstrating the
% LLNCS macro package for Springer Computer Science proceedings;
% Version 2.21 of 2022/01/12
%
\documentclass[runningheads]{llncs}
%
\usepackage[T1]{fontenc}
\usepackage[spanish]{babel}
% T1 fonts will be used to generate the final print and online PDFs,
% so please use T1 fonts in your manuscript whenever possible.
% Other font encondings may result in incorrect characters.
%
\usepackage{graphicx}
% Used for displaying a sample figure. If possible, figure files should
% be included in EPS format.
%
% If you use the hyperref package, please uncomment the following two lines
% to display URLs in blue roman font according to Springer's eBook style:
%\usepackage{color}
%\renewcommand\UrlFont{\color{blue}\rmfamily}
%\urlstyle{rm}
%
\begin{document}
%
\title{Evaluación de IPFS y Blockchain: Modelos de Almacenamiento y Distribución de Datos}
\titlerunning{Evaluación de IPFS y Blockchain}
%\titlerunning{Abbreviated paper title}
% If the paper title is too long for the running head, you can set
% an abbreviated paper title here
%
\author{Guillermo García Rodríguez\inst{1}}
%
\authorrunning{G. García Rodríguez}
% First names are abbreviated in the running head.
% If there are more than two authors, 'et al.' is used.
%
\institute{Universitat Politècnica de València \\
\email{guillegr1999@gmail.com}}
%
\maketitle              % typeset the header of the contribution
\makeatletter
\renewcommand{\l@title}[2]{}
\renewcommand{\l@author}[2]{}
\makeatother
%
\begin{abstract}
Este trabajo presenta una evaluación comparativa entre IPFS y Blockchain,
centrada en sus modelos de almacenamiento y distribución de datos.

\keywords{IPFS \and Blockchain \and Sistemas distribuidos \and Almacenamiento descentralizado \and Distribución de datos}
\end{abstract}

%
\setcounter{tocdepth}{2}
\tableofcontents
\cleardoublepage
%
\section{Introducción}

\section{Estado del arte}

\section{Arquitectura de IPFS y Blockchain}

\section{Integración de Inteligencia Artificial}

\section{Evaluación experimental}

\section{Resultados y discusión}

\section{Conclusiones}



\begin{thebibliography}{8}
\bibitem{ref_article1}
Author, F.: Article title. Journal \textbf{2}(5), 99--110 (2016)

\end{thebibliography}
\end{document}
